\section{Implementation on Zoo Network}

\label{sec:implementation}
%----------------------------------------------------------------------

\subsection{Multi-Chain Mapping}

The Memory Capsule architecture maps to \Zoo{} Network's specialized chains:

\begin{table}[h]
\centering
\begin{tabular}{lll}
\toprule
\textbf{Component} & \textbf{Chain} & \textbf{Responsibility} \\
\midrule
Capsule identity, anchors & \IChain{} & DID binding, head hash anchoring \\
Vault encryption keys & \KChain{} & Key derivation, rotation, escrow \\
Threshold reveal & \TChain{} & Share distribution, quorum decryption \\
FHE inference & \AChain{} & Homomorphic computation on encrypted state \\
Governance (GC policy) & \GChain{} & Community votes on default parameters \\
\bottomrule
\end{tabular}
\caption{Memory Capsule component mapping to \Zoo{} chains}
\label{tab:chain_mapping}
\end{table}

\subsection{A-Chain Integration}

\AChain{} provides the inference runtime for \FHE{}-encrypted memory queries. When an agent issues an encrypted recall:

\begin{enumerate}
    \item The encrypted query is submitted to \AChain{} as an inference request.
    \item \AChain{} validators execute the \CKKS{} similarity search across the capsule's encrypted embeddings.
    \item The encrypted result (top-$k$ indices) is returned to the agent.
    \item If threshold decryption is required, the result is forwarded to \TChain{}.
\end{enumerate}

\subsection{I-Chain Identity Binding}

Each capsule is bound to exactly one \DID{} on \IChain{}. The binding is enforced at the protocol level:

\begin{itemize}
    \item Capsule creation requires a valid \IChain{} \DID{} transaction.
    \item All operations on the capsule require a signature from the bound \DID{}.
    \item Ownership transfer is an \IChain{} transaction that atomically updates the \DID{} binding and re-encrypts the vault key via \KChain{}.
\end{itemize}

\subsection{K-Chain Key Hierarchy}

The key hierarchy for a capsule:

\begin{align}
k_{\textit{root}} &\leftarrow \text{KChain.DeriveRoot}(\textit{did}) \\
k_{\textit{vault}} &\leftarrow \text{HKDF}(k_{\textit{root}}, \text{``vault''}) \\
k_{\textit{entry}}(i) &\leftarrow \text{HKDF}(k_{\textit{vault}}, \text{``entry''} \| i) \\
k_{\textit{embedding}}(i) &\leftarrow \text{HKDF}(k_{\textit{vault}}, \text{``embed''} \| i)
\end{align}

Key rotation is handled by \KChain{}: a rotation event re-encrypts the vault key under a new root, without re-encrypting individual entries (which use derived keys from the vault key).

\subsection{T-Chain Threshold Protocol}

For threshold-protected memories, \TChain{} implements the following protocol:

\begin{enumerate}
    \item \textbf{Setup.} During capsule creation, the vault key $k_{\textit{vault}}$ is split into $n$ shares via Feldman VSS and distributed to \TChain{} validators.
    \item \textbf{Access.} The agent submits a decryption request with a valid capability token to \TChain{}.
    \item \textbf{Quorum.} Each validator independently verifies the capability token against \IChain{} and, if valid, contributes its share.
    \item \textbf{Reconstruction.} Once $t$ shares are collected, the key is reconstructed in a secure enclave, used for decryption, and immediately discarded.
\end{enumerate}

%----------------------------------------------------------------------
